\subsection{T1 Survived Local Affine Subtraction}

\begin{figure*}[t]
\centering
\includegraphics[width=\textwidth]{figures/Fig2_final.pdf}
\caption{T1 survival and full-pipeline calibration. The all-observation panel
retains both survivor and failure observations, the support matrix shows the
frozen two-scale decision structure, the omnibus-null panel shows the
distribution of both-scale support across $B=1000$ pseudo-event replicates,
and the detrending panel gives the main preprocessing boundary.}
\label{fig:t1_survival}
\end{figure*}

The first result was that local affine deformation did not fully absorb the
T1 observable. Under the two-scale consistency requirement, T1 survived in 14
of the 19 observations. Under the more permissive any-scale criterion, T1
survived in 15 of the 19 observations (Fig.~\ref{fig:t1_survival}). The
two-scale number is the main anchor because it shows that the survival
decision was not limited to a single original neighborhood-scale setting.

The all-observation support was also checked against a completed
full-pipeline pseudo-event calibration. In the observed data,
$N_{\mathrm{both}}=14$; across $B=1000$ null replicates, the null
distribution had mean 4.38 and q95 7, and no null replicate reached the
observed count of 14 (maximum 12; plus-one empirical
$p_{\mathrm{emp}}\approx0.001$). The any-scale count was supportive but less
specific, and is treated as secondary.

The detrending challenge preserved majority-level support but exposed an
important boundary. The frozen centered one-second detrending reproduced the
main $14/19$ both-scale result. A past-only one-second detrending reduced the
stricter both-scale count to $11/19$, while the no-rolling-detrending robust-z
variant gave $13/19$. Thus, the T1 survival result was not erased by the
challenge. The residual remained detectable at majority level under both
alternative preprocessing schemes, while the precise $14/19$ both-scale count
was specific to the frozen centered-detrending pipeline.

Local affine conditioning quality control did not indicate that the result was
driven by unstable local affine fits. Across all 38 observation-scale
combinations, all combinations passed the predefined conditioning screen, all
19 observations passed both $k$ values, the median condition number was 2.37,
the largest q95 condition number across combinations was 6.28, and no sampled
fit had condition number greater than 100.

The all-observation display is important for interpretation. Several
observations were fragile or failed under the predefined criteria. Thus, the
result supports common survival after local affine subtraction, not a
universal pattern across all observations.
